\documentclass[10pt,aspectratio=169]{beamer}
% Preámbulo modular para presentaciones Beamer (Modelación Estadística)
% Diseño y paleta institucional del Tecnológico de Monterrey

\documentclass[aspectratio=169,xcolor={dvipsnames,table}]{beamer}

\usepackage[utf8]{inputenc}
\usepackage[spanish,mexico]{babel}
\usepackage{amsmath,amssymb,amsfonts,mathtools}
\usepackage{graphicx}
\usepackage{listings}
\usepackage{booktabs}
\usepackage{tikz}
\usetikzlibrary{matrix,arrows,decorations.pathmorphing,shapes.geometric,positioning}

% Paleta Institucional Tecnológico de Monterrey
\definecolor{TecAzul}{HTML}{0039A6}       % Azul TEC: color primario institucional
\definecolor{TecAzulOscuro}{HTML}{1A2E51} % Azul marino oscuro
\definecolor{TecRojo}{HTML}{EC2661}       % Rojo de acento (paleta miTec)
\definecolor{TecGris}{HTML}{646464}       % Gris neutro para comentarios y subtítulos
\definecolor{TecFondoBloque}{HTML}{F4F6F9}% Fondo suave para bloques

% Configuración de Tema Beamer
\usetheme{Boadilla}
\usecolortheme[named=TecAzulOscuro]{structure}
\setbeamercolor{alerted text}{fg=TecRojo}
\setbeamercolor{palette primary}{bg=TecAzulOscuro,fg=white}
\setbeamercolor{palette secondary}{bg=TecAzul,fg=white}
\setbeamercolor{palette tertiary}{bg=TecAzulOscuro!80!black,fg=white}
\setbeamercolor{title}{fg=TecAzulOscuro,bg=TecFondoBloque}
\setbeamercolor{frametitle}{fg=TecAzulOscuro,bg=TecFondoBloque}

% Bloques personalizados elegantes
\setbeamercolor{block title}{bg=TecAzulOscuro,fg=white}
\setbeamercolor{block body}{bg=TecFondoBloque,fg=black}
\setbeamercolor{block title alerted}{bg=TecRojo,fg=white}
\setbeamercolor{block body alerted}{bg=TecRojo!8,fg=black}
\setbeamercolor{block title example}{bg=TecAzul,fg=white}
\setbeamercolor{block body example}{bg=TecAzul!8,fg=black}

% Redondear bordes y añadir sombra sutil
\setbeamertemplate{blocks}[rounded][shadow=true]
\setbeamertemplate{navigation symbols}{}

% Pie de página limpio con número de diapositiva
\setbeamertemplate{footline}{
  \leavevmode%
  \hbox{%
  \begin{beamercolorbox}[wd=.7\paperwidth,ht=2.25ex,dp=1ex,left]{author in head/foot}%
    \hspace*{2ex}\insertshortauthor~~(\insertshortinstitute) \hfill \insertshorttitle
  \end{beamercolorbox}%
  \begin{beamercolorbox}[wd=.3\paperwidth,ht=2.25ex,dp=1ex,right]{date in head/foot}%
    \insertshortdate{}\hspace*{2em}
    \insertframenumber{} / \inserttotalframenumber\hspace*{2ex}
  \end{beamercolorbox}}%
  \vskip0pt%
}

% Configuración de Listings de Python para visualización pedagógica de código
\lstset{
    language=Python,
    basicstyle=\ttfamily\footnotesize,
    keywordstyle=\color{TecAzulOscuro}\bfseries,
    stringstyle=\color{TecRojo},
    commentstyle=\color{TecGris}\itshape,
    showstringspaces=false,
    breaklines=true,
    frame=lines,
    rulecolor=\color{TecAzulOscuro!30},
    backgroundcolor=\color{TecFondoBloque},
    aboveskip=0.5em,
    belowskip=0.5em,
    literate={á}{{\'a}}1 {é}{{\'e}}1 {í}{{\'i}}1 {ó}{{\'o}}1 {ú}{{\'u}}1
             {Á}{{\'A}}1 {É}{{\'E}}1 {Í}{{\'I}}1 {Ó}{{\'O}}1 {Ú}{{\'U}}1
             {ñ}{{\~n}}1 {Ñ}{{\~N}}1 {¿}{{?`}}1 {¡}{{!`}}1
}

% Metadatos institucionales por defecto
\title[Modelación Estadística]{Modelación Estadística para Ciencia de Datos e Ingeniería}
\author[J. Castillo Colmenares]{Juliho David Castillo Colmenares}
\institute[Tec de Monterrey]{Tecnológico de Monterrey \\ Escuela de Ingeniería y Ciencias}
\date{\today}
% Comandos y macros estadísticas compartidas para las presentaciones Beamer (Metropolis)

% Conjuntos numéricos básicos
\newcommand{\R}{\mathbb{R}}
\newcommand{\N}{\mathbb{N}}
\newcommand{\Q}{\mathbb{Q}}
\newcommand{\Z}{\mathbb{Z}}
\newcommand{\set}[1]{\left\{ #1 \right\}}
\newcommand{\abs}[1]{\left|#1\right|}
\newcommand{\norm}[1]{\left\| #1 \right\|}

% Comandos estadísticos y probabilísticos del curso
\newcommand{\Var}{\operatorname{Var}}
\newcommand{\cov}{\operatorname{Cov}}
\newcommand{\corr}{\rho}
\newcommand{\comb}[2]{\begin{pmatrix} #1 \\ #2 \end{pmatrix}}
\newcommand{\s}{\sigma}
\newcommand{\particion}{\mathfrak{P}}
\newcommand{\card}[1]{\#\left( #1 \right)}
\newcommand{\seq}[3]{\set{#1}_{#2}^{#3}}
\newcommand{\E}{\mathbb{E}}
\newcommand{\Prob}{\mathbb{P}}

% Entornos pedagógicos y cajas en español académico natural
\newenvironment{discusion}{%
  \begin{alertblock}{Pregunta de Discusión}%
}{%
  \end{alertblock}%
}

\newenvironment{reflexion}{%
  \begin{alertblock}{Reflexión para el Aula}%
}{%
  \end{alertblock}%
}

\newenvironment{paradoja}{%
  \begin{alertblock}{Paradoja de Exploración}%
}{%
  \end{alertblock}%
}

\newenvironment{motivacion}{%
  \begin{exampleblock}{Motivación: Ciencia de Datos en la Práctica}%
}{%
  \end{exampleblock}%
}

\newenvironment{retoclase}{%
  \begin{block}{Reto Rápido en Equipo}%
}{%
  \end{block}%
}

\title[Pruebas de homogeneidad]{Sección 07.11: Pruebas de homogeneidad}\subtitle{Una variable categórica en varios grupos}
\author[J. Castillo Colmenares]{Juliho Castillo Colmenares}\institute{Tecnológico de Monterrey}\date{}
\begin{document}
\begin{frame}[plain]\titlepage\end{frame}
\begin{frame}{Hoja de ruta}\small\begin{enumerate}\item Diseño y concepto.\item Tabla de frecuencias.\item Estadística y decisión.\end{enumerate}\end{frame}
\begin{frame}{Fuente y alcance}\small La prueba compara la distribución de una única variable categórica entre varios subgrupos independientes.\begin{block}{Pregunta guía}¿Comparten los grupos la misma distribución categórica?\end{block}\end{frame}
\begin{frame}{Homogeneidad e independencia}\small La fórmula chi-cuadrada coincide con la de independencia, pero el diseño cambia: aquí se fijan muestras por población y se compara una variable.\end{frame}
\begin{frame}{Diseño muestral}\small Se extraen $c$ muestras independientes de tamaños $n_1,\ldots,n_c$ y se observan $r$ categorías en cada muestra.\end{frame}
\begin{frame}{Hipótesis}\small\[H_0:\text{las }c\text{ poblaciones tienen la misma distribución};\qquad H_a:\text{alguna difiere}.\]\end{frame}
\begin{frame}{Frecuencias esperadas}\small Para una tabla $r\times c$,\[E_{ij}=\frac{(\text{total fila }i)(\text{total columna }j)}{\text{total general}}.\]\end{frame}
\begin{frame}{Estadística chi-cuadrada}\small\[\chi^2=\sum_{i=1}^{r}\sum_{j=1}^{c}\frac{(O_{ij}-E_{ij})^2}{E_{ij}}.\]\end{frame}
\begin{frame}{Grados de libertad}\small La estadística se compara con $\chi^2_{(r-1)(c-1)}$. Las restricciones de los totales reducen la información libre.\end{frame}
\begin{frame}{Condiciones}\small Las observaciones deben ser independientes y las frecuencias esperadas suficientemente grandes para justificar la aproximación chi-cuadrada.\end{frame}
\begin{frame}{Procedimiento I}\small\begin{enumerate}\item Definir grupos y categorías.\item Construir la tabla observada.\item Calcular frecuencias esperadas.\end{enumerate}\end{frame}
\begin{frame}{Procedimiento II}\small\begin{enumerate}\setcounter{enumi}{3}\item Calcular $\chi^2$.\item Obtener valor crítico o $p$-value.\item Interpretar la distribución, no solo una celda.\end{enumerate}\end{frame}
\begin{frame}{Lectura de la estadística}\small Celdas con $O_{ij}$ muy distinto de $E_{ij}$ aportan más a $\chi^2$. La suma agrega discrepancias de toda la tabla.\end{frame}
\begin{frame}{Aplicación: tabla}\small Organice las frecuencias de una variable categórica en filas y los grupos independientes en columnas. Los totales de columna reflejan el diseño.\end{frame}
\begin{frame}{Aplicación: esperadas}\small Calcule cada $E_{ij}$ con los totales marginales. Compruebe que las esperadas respeten los totales de filas y columnas.\end{frame}
\begin{frame}{Aplicación: contraste}\small Evalúe $\chi^2$ con $gl=(r-1)(c-1)$ y nivel $\alpha$.\end{frame}
\begin{frame}{Aplicación: decisión}\small Si el valor $p$ es pequeño, hay evidencia de que no todas las distribuciones son iguales.\end{frame}
\begin{frame}{Qué no identifica la prueba}\small Un rechazo global no indica qué grupos o categorías explican la diferencia. Se requieren comparaciones posteriores y una lectura sustantiva.\end{frame}
\begin{frame}{Errores frecuentes}\small\begin{itemize}\item Llamar independencia al diseño de homogeneidad.\item Ignorar esperadas pequeñas.\item Confundir rechazo global con diagnóstico de una celda.\end{itemize}\end{frame}
\begin{frame}{Plantilla de reporte}\small\begin{block}{Reporte} ``La tabla $r\times c$ produjo $\chi^2=\cdots$, $gl=\cdots$, $p=\cdots$. Al nivel $\alpha=\cdots$, [rechazamos/no rechazamos] la homogeneidad.''\end{block}\end{frame}
\begin{frame}{Síntesis}\small\begin{itemize}\item Homogeneidad compara poblaciones mediante una variable categórica.\item La estadística usa observadas y esperadas.\item El rechazo global requiere análisis posterior.\end{itemize}\end{frame}
\begin{frame}{Conexión con el capítulo}\small La prueba es el marco general para el contraste de varias proporciones que sigue.\end{frame}
\end{document}
